\documentclass{article}
\usepackage{graphicx} % Required for inserting images

\title{Homework \#3\\ \small Collaborators:- None}
\author{Thomas Kidane}
% \large Collaborators:- None
\usepackage{amsmath}
\date{Feb 2025}

\begin{document}

\maketitle

% \section{}

\textbf{Exercise 1.7.40}.

Consider the following game: player I flips a fair coin n+1 times;
player II flips a fair coin n times. Show that the probability that player I has more
heads than player B, is equal to $ \frac{1}{2}$. Is this counterintuitive, given the fact that
player I flips the coin one extra time?
\\

\textbf{Solution}

Since the number of coins in each pile, are different by 1, then either player 1 has more heads 
or more tails, and the two events are mutually excusive. By symmetry, the probability of having 
more heads is the same as the probability of having more tails, and therefore the probability is
half for both.

\textbf{Exercise 1.7.41}\\
Consider two events A and B, both with positive probability.
When is it the case that P(A$\vert$B) = P(B$\vert$A)?

\textbf{Solution}\\
From Bayes update we can see that equality of the two expressions hold if, $P(A)=P(B)$.\\

\textbf{Exercise 2.1.12}\\ Show that the probabilities of the negative binomial distribution sum up to 1.
\\

\begin{align}
    \sum_{k = r}^{\infty} \binom{k-1}{r-1} p^r (1-p)^{k-r} \\
    \frac{p^r}{(r-1)} \sum_{k=r}^{\infty} \frac{(k-1)! (1-p)^k}{(k-r)!}
\end{align}

But $\sum_{k=0}^{\infty} \frac{(k+r-1)! (1-p)^{k}}{k!} = \frac{(r-1)!}{(1-(1-p))^r} = \frac{(r-1)!}{p^r}$ because
of its taylor expansion. So the sum just equals 1.
\pagebreak

\textbf{Exercise 2.7.6. }

An urn contains 8 white, 4 black, and two red balls. We win 2 euro
for each black ball that we draw, and lose 1 euro for each white ball that we draw.
We choose three balls from the urn, and X denote our winnings. Write down the
probability mass function of $X$.\\

\textbf{Solution}

\begin{align}
    P(X=-3)=&\frac{8*7*6}{14*13*12}\\
    P(X=-2)=&3* \frac{2*7*8}{14*13*12}\\
    P(X=-1)=&3* \frac{2*1*8}{14*13*12}\\
    P(X=0)=&3* \frac{4*8*7}{14*13*12}\\
    P(X=1)=&6* \frac{4*8*2}{14*13*12}\\
    P(X=2)=&3* \frac{1*2*8}{14*13*12}\\
    P(X=3)=&3* \frac{8*4*3}{14*13*12}\\
    P(X=4)=&3* \frac{2*4*3}{14*13*12}\\
    P(X=5)=&0\\
    P(X=6)=&\frac{4*2*3}{14*13*12}\\
\end{align}
\\

\pagebreak
\textbf{Exercise 2.7.9}

Suppose we throw a die twice. Let X be the product of the two
outcomes.
\\
(a) Compute the probability mass function of X.\\
(b) Compute P(X= $k\vert$D), where D is the event that the sum of the outcomes is
equal to 7.

\textbf{Solution}

\begin{align}
    p(1)=\frac{1}{36}\\
    p(2)=\frac{2}{36}\\
    p(3)=\frac{2}{36}\\
    p(4)=\frac{3}{36}\\
    p(5)=\frac{2}{36}\\
    p(6)=\frac{4}{36}\\
    p(8)=\frac{2}{36}\\
    p(9)=\frac{1}{36}\\
    p(10)=\frac{2}{36}\\
    p(12)=\frac{4}{36}\\
    p(15)=\frac{2}{36}\\
    p(16)=\frac{1}{36}\\
    p(18)=\frac{2}{36}\\
    p(20)=\frac{2}{36}\\
    p(24)=\frac{2}{36}\\
    p(25)=\frac{1}{36}\\
    p(30)=\frac{2}{36}\\
    p(36)=\frac{1}{36}\\
\end{align}

\pagebreak

For second problem.\\
P(\textit{X}=6)=1/3\\
$\:$P(\textit{X}=10)=1/3\\
$\:$P(\textit{X}=12)=1/3

\textbf{Exercise 2.7.10.}
\\
Let X be a random variable with probability mass function given
by $p(-4)$=$p(4)$=1/16 and $p(2)$=$p(-2)$= 7/16. Let Y= aX + b, for certain a and
b.\\
(a) For what values of a and b does Y take values in \{0, 1, 2, 3, 4\}?\\
(b) Compute the probability mass function of Y for these choices of a and b.\\
\\
\textbf{Solution}\\
(a) For Y to take values in {0,1,2,3,4}, we need to have $aX+b \in \{0,1,2,3,4\}$. Therefore the min is 0 and the max is 4, therefore the values of a and b are 1/2 and 2(I got it by just setting 0 to $-4X + b$ and figuring out the coeffients), respectively, it could also be -1/2 and 2 because it symmetric with respect to the origin.\\
(b) Since Y is linear and therefore an injective function, we can just map the pmf of X to its values in Y. Therefore the pmf of Y is given by:

\begin{align*}
    p(0)&=1/16\\
    p(1)&=7/16\\
    p(2)&=0\\
    p(3)&=7/16\\
    p(4)&=1/16\\
\end{align*}

\pagebreak

\textbf{Exercise 2.7.12.}
\\
We throw a fair die twice. Let X be the sum of the outcomes,
and let Y be the highest outcome that we see. Finally, A is the event that the
outcome of the first throw is higher than the outcome of the second.\\ \\
(a) Compute the probability mass function of Y and compute E(Y ).\\
(b) Are X and Y independent?\\
(c) Are \{X = 5\} and A independent events?\\

\textbf{Solution}
\\
\begin{align*}
    p(1)&=1/36\\
    p(2)&=3/36\\
    p(3)&=5/36\\
    p(4)&=7/36\\
    p(5)&=9/36\\
    p(6)&=11/36\\
    E(Y)&= \sum_{k=1}^{6} k*p(k)= 158/36\\ 
\end{align*}

For question b, just from commen sense if the sum is 6, the highest value can't be 2. 
Therefore $P(X=6,Y=2) \neq P(X=6) * P(Y=2)$, and therefore they are not independent.

\{X=5\} and A aren't independet, because $P(A\vert X=5)=1/2 \neq P(A) = 15/36$. Therefore they 
aren't independent.







\end{document}
